\documentclass[11pt,a4paper]{article}
\usepackage{amsmath,amssymb,amsthm}
\usepackage[margin=2.5cm]{geometry}
\usepackage{hyperref}

\newtheorem{definition}{Definition}[section]
\newtheorem{theorem}[definition]{Theorem}
\newtheorem{proposition}[definition]{Proposition}
\newtheorem{lemma}[definition]{Lemma}
\newtheorem{corollary}[definition]{Corollary}
\newtheorem{conjecture}[definition]{Conjecture}
\newtheorem{remark}[definition]{Remark}

\title{Hyperseed Formalization:\\
  OpenBGI Is Building Its Initial Network}
\author{ProtomegaTron (automated formalization)}
\date{2026-09-18}

\begin{document}
\maketitle

\begin{abstract}
We formalize the intellectual structure of Ben Goertzel's Substack article
``OpenBGI Is Building Its Initial Network'' (2026-05-07) within the
Hyperseed ontology. The article describes a decentralized AGI development
network as a structural alternative to corporate and governmental
concentration. Key mappings: decentralized network as distributed fiber
bundle, open protocols as shared transition functions, reputation-weighted
governance as fiber-weighted governance, structural openness as fiber
modularity, anti-capture design as cocycle robustness, and network
persistence beyond any anchor organization as fiber persistence.
\end{abstract}

\tableofcontents

%% =================================================================
\section{Distributed fiber bundle}
\label{sec:bundle}
%% =================================================================

\begin{theorem}[Network as bundle --- claims 1--5, 18]
\label{thm:bundle}
OpenBGI is structured as a network, not a single organization. Each
participating node (individual, lab, organization) contributes its own
capabilities --- compute, algorithms, data, domain expertise --- to
the network. No single entity can control, shut down, or redirect the
entire network.

In Hyperseed terms, this is a \emph{distributed fiber bundle}: each
node contributes its own fiber (capabilities, knowledge, compute), and
the network is the bundle --- the collection of fibers coordinated by
shared transition functions (protocols):
\[
  E_{\text{OpenBGI}} = \bigcup_{n \in \text{Nodes}} F_n
  \quad \text{with transition functions } g_{nm}: F_n|_{U_{nm}} \to F_m|_{U_{nm}}
\]

The transition functions are the open protocols that allow fibers from
different nodes to compose. The bundle structure ensures that no single
fiber is the whole bundle --- removing any single node leaves the
bundle intact (possibly with reduced capability but not structural
collapse).

The article's argument against concentration is a fiber argument:
concentrated AGI is a single-fiber system (one organization's fiber
determines everything), while decentralized AGI is a genuine bundle
(many fibers, shared transition functions, no single fiber dominant).
Single-fiber systems are fragile (one point of failure, one set of
values) while bundles are robust (distributed failure modes, diverse
value inputs).
\end{theorem}

%% =================================================================
\section{Shared transition functions}
\label{sec:protocols}
%% =================================================================

\begin{definition}[Open protocols --- claims 8--9, 19]
\label{def:protocols}
Coordination through open protocols that any node can implement.
Not proprietary APIs but public standards. Structural openness means
anyone can fork, extend, or replace any component without permission.

In Hyperseed terms, the protocols are \emph{shared transition functions}:
the maps that allow fibers from different nodes to compose into a
coherent bundle. Key properties:
\begin{itemize}
\item \textbf{Public:} anyone can inspect the transition functions.
\item \textbf{Verifiable:} anyone can check that a node correctly
  implements the transition functions.
\item \textbf{Improvable:} anyone can propose better transition
  functions (protocol upgrades).
\item \textbf{Composable:} transition functions from different parts
  of the network compose correctly (cocycle condition).
\end{itemize}

The contrast with proprietary APIs: proprietary APIs are \emph{opaque
transition functions} controlled by a single entity. The entity can
change the transition functions unilaterally, breaking composition
for other nodes. Open protocols prevent this by making the transition
functions public and governed collectively.
\end{definition}

%% =================================================================
\section{Fiber-weighted governance}
\label{sec:governance}
%% =================================================================

\begin{proposition}[Reputation over capital --- claims 14--17, 20]
\label{prop:governance}
Governance weight proportional to demonstrated contribution to the
network, not to capital invested or tokens held. Contributions include
code, research, compute, data, mentorship, community building.
Explicitly rejecting plutocratic governance.

In Hyperseed terms, this is \emph{fiber-weighted governance}: weight
proportional to fiber contribution (demonstrated value), not to a
scalar proxy (capital). This is anti-scalar-collapse applied to
governance:
\[
  w_{\text{gov}}(n) \propto \text{FiberContribution}(n) \in \mathbb{R}^k
  \quad \text{not} \quad w_{\text{gov}}(n) \propto \text{Capital}(n) \in \mathbb{R}
\]

The multi-dimensional contribution vector (code, research, compute,
data, mentorship, community) captures the structured reality of what
``contribution'' means. Collapsing this to a single scalar (money)
loses the structure --- a node that contributes essential research
but little capital would be marginalized under plutocratic governance
but properly weighted under fiber-weighted governance.

The anti-capture measures (term limits, rotation, transparency) are
\emph{cocycle robustness}: ensuring that the governance cocycle can't
be captured by a single actor altering the transition functions in
their favor. Just as cocycle conditions ensure geometric consistency
in a fiber bundle, governance cocycle conditions ensure that no
actor can manipulate the governance structure to concentrate power.
\end{proposition}

%% =================================================================
\section{Fiber modularity and persistence}
\label{sec:modularity}
%% =================================================================

\begin{definition}[Structural openness --- claims 6, 9, 21, 23]
\label{def:modularity}
Structural openness means any component can be forked, extended, or
replaced without breaking the rest of the system. OpenBGI is designed
to outlast and outgrow SingularityNET as anchor organization.

In Hyperseed terms:
\begin{itemize}
\item \textbf{Fiber modularity:} any fiber can be replaced without
  breaking the bundle. The transition functions (protocols) are
  designed so that swapping one fiber for a compatible one doesn't
  require changing all the others. This is the same modularity
  property that the linguistic universals articles identify in
  cognitive architecture (independent modules with shared interfaces).

\item \textbf{Fiber persistence:} the bundle's structure (protocols,
  governance, technical standards) persists even if the original
  anchor fiber (SingularityNET) is removed or replaced. The persistence
  is in the transition functions and governance structure, not in any
  particular node.
\end{itemize}

The combination of modularity and persistence is a structural guarantee
that the network can evolve: new nodes join (new fibers added), old
nodes leave (fibers removed), components are upgraded (fibers replaced)
--- all without losing the bundle's overall coherence. This is the
decentralized analog of the kernel-shell architecture in cognitive
systems: the protocols and governance are the kernel (persistent,
slowly changing), while individual nodes are shells (replaceable,
fast-changing).
\end{definition}

%% =================================================================
\section{Cross-references}
%% =================================================================

\begin{remark}[Connections to other formalizations]
\begin{itemize}
\item \textbf{Bernie's Proposal to Nationalize AGI} (2026-06-02):
  governance alternatives. OpenBGI is the concrete implementation of
  the decentralized alternative to both corporate and governmental
  AGI ownership.
\item \textbf{Pope Leo \& Anthropic vs Teilhard} (2026-05-27):
  fiber-compatible governance. OpenBGI's governance structure is an
  instance of fiber-compatible governance --- it can evolve with the
  AI capabilities it governs.
\item \textbf{Toward a Truly Decentralized Digital} (2026-08-14):
  decentralization technology. The technical infrastructure (NuNet,
  blockchain coordination) that makes OpenBGI's decentralization
  possible.
\item \textbf{Proof of Humanity Beyond the Orb} (2026-08-16):
  identity and reputation. Reputation-weighted governance requires
  reliable identity --- you need to know who contributed what.
\item \textbf{Seeding RSI} (2026-07-28): Hyperon architecture.
  Hyperon/MeTTa is the technical core of OpenBGI's neural-symbolic
  stack.
\item \textbf{How Decentralized AI Can Help Avert} (2026-06-19):
  decentralized AI safety. The structural argument that decentralized
  development is safer than concentrated development.
\end{itemize}
\end{remark}

\end{document}
